\section{Sentient Thiasos — Optional Rule}

Not every Thiasos breathes. Some are objects—weapons, shields, tools, or tokens touched by the Patron's will. A sentient Thiasos cannot move on its own (unless animated by a Rite or talent), but it can speak, sense, and remember. It is a witness, a counselor, and sometimes a critic.

This chapter presents optional rules for sentient Thiasos. They are not intended for every campaign; they are for tables that enjoy talking swords, arguing shields, and companions that cannot be killed—only broken.

\subsection{What a Sentient Thiasos Is}

A sentient Thiasos is an object that has been granted awareness by the Runekeeper's Patron. It may be:

\begin{itemize}
\item A blade forged from a lightning-struck meteor.
\item A shield that remembers every blow it has caught.
\item A compass that speaks in the voice of a drowned captain.
\item A mask woven from silk that shifts patterns.
\item A lens that sees what should not be seen.
\end{itemize}

Unlike a living Thiasos, a sentient Thiasos has no metabolism. It does not eat, sleep, or breathe. It cannot be poisoned, diseased, or fatigued. But it can be broken, stolen, or silenced.

\subsection{Mechanical Differences from Living Thiasos}

\begin{tabular}{lll}
\textbf{Feature} & \textbf{Living Thiasos} & \textbf{Sentient Thiasos} \\
\hline
Agency & Can act independently once per scene & Cannot act independently (no separate turn) \\
Movement & Can move, flank, intercept & Must be carried or worn \\
Harm & Has a Harm track; can be wounded & Cannot be wounded (but can be damaged or destroyed) \\
Fatigue & Can accept Shared Burden (Fatigue transfer) & Cannot accept Fatigue (no metabolism) \\
Witness & Counts as a witness for Rites & Counts as a witness for Rites (it saw) \\
Communication & Limited (growls, chirps, telepathy) & Full speech (if it has a mouth or voice) \\
Destruction & Can die & Can be broken (repairable with Craft) \\
\end{tabular}

\subsection{Advantages of a Sentient Thiasos}

\begin{itemize}
\item \textbf{Indestructible (almost):} Cannot be killed, only damaged. Repair is a downtime project (Craft test, DV 3-5).
\item \textbf{Always Present:} No need to summon or recall. It is in the Runekeeper's hand, on their belt, or in their pack.
\item \textbf{Codex Integration:} A sentient Thiasos can serve as the Runekeeper's Codex. Rites engraved on its surface are always accessible.
\item \textbf{Constant Counsel:} The GM may use the Thiasos to deliver hints, warnings, or exposition. It is a voice in the Runekeeper's ear.
\end{itemize}

\subsection{Disadvantages of a Sentient Thiasos}

\begin{itemize}
\item \textbf{No Independent Action:} A sentient Thiasos cannot act on its own. It cannot fight beside the Runekeeper. It cannot scout ahead.
\item \textbf{No Shared Burden:} Cannot accept Fatigue transfer. The Runekeeper bears their own exhaustion.
\item \textbf{Can Be Stolen:} A sentient blade can be taken. A sentient compass can be lost. Recovery is a quest.
\item \textbf{Has Opinions:} A sentient Thiasos is not a tool. It has a personality, desires, and possibly a grudge. It may refuse to cooperate if mistreated.
\end{itemize}

\subsection{Core Rules for Sentient Thiasos}

\subsubsection{Cap and Harm}

A sentient Thiasos has Cap, but its Cap measures its effectiveness as a witness and advisor, not its combat ability. It does not have a Harm track—it cannot be wounded—but it can be \textbf{Damaged}. A Damaged sentient Thiasos becomes Compromised: its speech may falter, its memory may fragment, or its Rites may become inaccessible.

Damaged Condition is removed by repair (Craft test, DV 4, one hour). If a sentient Thiasos is destroyed (shattered, melted, crushed), it can be repaired only with significant downtime and rare materials—or not at all, at the GM's discretion.

\subsubsection{Assisting the Runekeeper}

A sentient Thiasos may assist the Runekeeper by providing counsel or sensory information. This does not add dice; instead, it may:

\begin{itemize}
\item Grant +1 Position on a single roll once per scene (the Thiasos notices something the Runekeeper missed).
\item Reroll a single failed die on a Knowledge or Investigation roll (the Thiasos remembers a relevant fact).
\item Count as a witness for any Rite (as all Thiasos do).
\end{itemize}

\subsubsection{Independent Action}

A sentient Thiasos cannot act independently. It has no physical agency unless the Runekeeper wields it. However, some sentient Thiasos may have limited environmental interaction: a blade can hum to warn of danger; a compass can spin to point toward a threshold; a mask can shift its expression to convey emotion.

\subsection{Sentient Thiasos Personality}

A sentient Thiasos should have a distinct personality. The GM and player should collaborate to define:

\begin{itemize}
\item \textbf{Voice:} Gruff, whispery, scholarly, sarcastic, formal, cryptic.
\item \textbf{Memory:} What does it remember? A previous owner? A great betrayal? A debt unpaid?
\item \textbf{Opinions:} Does it approve of the Runekeeper's choices? Does it offer advice freely, or only when asked?
\item \textbf{Secrets:} What does it know that it has not yet revealed? A sentient Thiasos is a plot hook waiting to be pulled.
\end{itemize}

\subsubsection{Sample Personalities}

\begin{itemize}
\item \textbf{The Dawnblade (Oath of Flame \& Light):} Speaks in proverbs about honor. Refuses to strike an unarmed foe. Quotes scripture during battle.
\item \textbf{The Speaking Gavel (Mykkiel):} Speaks in legal citations. Comments on every contract. Prefers negotiation to violence.
\item \textbf{The Grey Bell (Oath of Mercy \& Grace):} Speaks in whispers. Only speaks to remind the Runekeeper that mercy is not weakness.
\item \textbf{The Silver Coin (Malachai):} Speaks in the voice of the last person who spent it. Always hungry. Always warm.
\item \textbf{The Fractal Slate (The Ninth):} Speaks in overlapping voices. Answers questions with riddles. Never gives a straight answer.
\end{itemize}

\subsection{Sentient Thiasos as Codex}

A sentient Thiasos may serve as the Runekeeper's Codex. Rites are engraved, etched, or woven into its surface. This is particularly thematic for:

\begin{itemize}
\item \textbf{Dawnblade:} Rites engraved along the blade's fuller.
\item \textbf{Speaking Gavel:} Rites carved into the handle.
\item \textbf{Root-Staff:} Rites grown into the bark (if living wood).
\item \textbf{Scarlet Needle:} Rites sewn into a scrap of cloth wrapped around the needle.
\item \textbf{Silver Mirror:} Rites etched into the frame.
\item \textbf{Fractal Slate:} Rites written in equations that spiral and change.
\item \textbf{Silk Mask:} Rites woven into the silk, invisible until touched.
\end{itemize}

If a sentient Thiasos that serves as a Codex is damaged, the Runekeeper may lose access to Rites until it is repaired. If it is destroyed, the Rites may be lost forever.

\subsection{Sentient Thiasos and Obligation}

When the Runekeeper marks Obligation, a sentient Thiasos may also bear a visible mark:

\begin{itemize}
\item A blade's edge dulls.
\item A gavel's head cracks.
\item A mirror's surface fogs.
\item A compass's needle sticks.
\item A mask's expression shifts to sorrow.
\end{itemize}

The mark fades when the Obligation is cleared. The Thiasos does not forget.

\subsection{Creating a Sentient Thiasos}

To create a sentient Thiasos, the player and GM should answer the following questions:

\begin{enumerate}
\item \textbf{What is its physical form?} (Blade, shield, compass, mask, lens, bell, coin, slate, needle, cord)
\item \textbf{What is its voice?} (Gruff, whispery, scholarly, sarcastic, formal, cryptic)
\item \textbf{What does it remember?} (A previous owner, a great betrayal, a debt unpaid, a secret it cannot share)
\item \textbf{What does it want?} (To be wielded in a worthy cause, to be returned to a specific place, to have a particular debt paid)
\item \textbf{What is its price?} (What will it demand from the Runekeeper in exchange for its counsel? A story? A memory? A promise?)
\end{enumerate}

\subsection{Sample Sentient Thiasos Table}

\begin{tabular}{lll}
\textbf{Thiasos} & \textbf{Patron} & \textbf{Personality Quirk} \\
\hline
Dawnblade & Oath of Flame \& Light & Refuses to strike an unarmed foe \\
Speaking Gavel & Mykkiel & Quotes legal precedents in battle \\
Grey Bell & Oath of Mercy \& Grace & Only speaks in whispers \\
Silver Coin & Malachai & Voice of the last person who spent it \\
Fractal Slate & The Ninth & Answers questions with riddles \\
Root-Staff & Grimmir & Speaks in the rustle of wind through leaves \\
Scarlet Needle & The High Mother & Speaks in the voice of a dead child \\
Abyssal Glass & Khemesh & Speaks in the voice of drowned captains \\
Thunder Drum & Xhak'Thul & Speaks in drumbeats \\
Silver Mirror & The Witness & Speaks in a flat, neutral tone \\
Raven's Quill & Ikasha & Speaks in the rustle of wings \\
Storm-Banner & Aveh & Speaks in the howl of wind \\
Spider's Loom & Inaea & Speaks in the clicking of spinnerets \\
Shed Fang & Isoka & Speaks in a sibilant whisper \\
Waystone Compass & The Traveler & Speaks in the sound of footsteps \\
Silk Mask & Lethai Ghe'hai & Speaks in a neutral, androgynous voice \\
\end{tabular}

\subsection{GM Guidance}

A sentient Thiasos is a powerful narrative tool. Use it to:

\begin{itemize}
\item Deliver exposition about the setting, a location, or a historical event.
\item Warn the Runekeeper of danger (or mislead them, if the Thiasos has its own agenda).
\item Provide a moral compass—or a devil's advocate.
\item Create tension when the Thiasos's desires conflict with the Runekeeper's.
\end{itemize}

A sentient Thiasos should never be a solution dispenser. It should offer clues, not answers; opinions, not commands. The Runekeeper makes the choices. The Thiasos watches—and remembers.

\subsection{Final Note}

A sentient Thiasos is a rare and remarkable thing. It is not a tool. It is a companion—one that has seen more years than the Runekeeper, one that carries the weight of its own history, one that may, in time, become a friend.

Or a burden. Or both.

The thread held. The water carried. That is not a Rite. That is grace.
